\section{Software architecture}
The five tools share one library (\texttt{crisprlib}) and one command-line
front end (\texttt{cli.py}); each tool package owns only its domain logic.

\subsection{Module map}
\texttt{crisprlib.featurize} holds every sequence encoding (one-hot,
dinucleotide one-hot, $k$-mer counts, GC content, SantaLucia nearest-neighbour
$\Delta G$, microhomology scoring, guide:target alignment encodings).
\texttt{crisprlib.models} holds GuideCNN, PairCNN, and SequenceGCN with their
graph construction. \texttt{crisprlib.train} holds the deterministic CPU
training loop (Adam, MSE, early stopping, fixed seeds) shared by every model
in the paper, so all training behaviour described in Methods exists exactly
once in the codebase. \texttt{crisprlib.benchmark} holds the evaluation
metrics (per-gene Spearman, regression metrics) used by every benchmark
script, preventing metric drift between tools.

\subsection{Tool packages}
\texttt{tools\_multiplex} (duplex screening, shared-PAM competition scoring,
pair ranking), \texttt{tools\_outcome} (inDelphi loading, event aggregation,
frameshift/insertion/deletion targets), \texttt{tools\_riskflag}
(microhomology density, secondary-structure and exon-boundary terms, locus
loader), \texttt{tools\_biosecurity} (homology pre-filter, audit-record
writer, Common Mechanism handoff), and \texttt{tools\_transfer} (FC+RES/V1
loaders with the contamination guard, divergence statistics, transfer
evaluator). The loaders are the single place where the
\texttt{drug} $\neq$ \texttt{nodrug} filter is enforced; benchmark scripts
cannot bypass it without editing library code, which is how the contamination
class from Section on negative results is now structurally prevented.

\subsection{Command-line surface}
\texttt{python cli.py <tool> <subcommand>} exposes every user-facing
function: \texttt{multiplex screen}, \texttt{outcome annotate},
\texttt{riskflag score}, \texttt{biosecurity prescreen}, \texttt{transfer
evaluate}. All emit JSON to stdout and write audit records where relevant;
there are no interactive prompts, so every tool composes in shell pipelines.

\subsection{Design decisions and their costs}
Three decisions were made for auditability over raw performance. (1) Pure
NumPy/PyTorch implementations with no CUDA paths: every number reproduces on
free-tier CPUs in minutes, at the cost of excluding very large screens.
(2) JSON-over-files artifacts: every result the paper cites is a committed
file, at the cost of a larger repository. (3) A single shared training loop:
cross-tool comparability, at the cost of forgoing per-model exotic
optimizers. We judge these costs acceptable for a benchmark-and-discovery
deliverable whose value depends on third parties being able to re-run it.

\subsection{Failure containment}
Each script validates its inputs against \texttt{data/PROVENANCE.md} hashes
where network data is involved, refuses to overwrite existing result files
unless told to, and exits non-zero on any NaN in model outputs. The 40-test
hermetic suite covers featurization invariants (one-hot row sums, graph
symmetry), loader guards (no \texttt{nodrug} rows in training frames),
metric identities, and end-to-end smoke runs of each tool on synthetic
fixtures; it needs no network and finishes in under a minute.
